\section{Building and Campus Ledger}

The ledger distinguishes construction, alteration, reported damage, and planned work. A date in the first column identifies the best-supported project period, not a claim that every component was completed that year.

\begin{longtable}{>{\bfseries\RaggedRight\arraybackslash}p{0.15\linewidth}>{\RaggedRight\arraybackslash}p{0.43\linewidth}>{\RaggedRight\arraybackslash}p{0.34\linewidth}}
\toprule
\textbf{Date} & \textbf{Building act or physical state} & \textbf{Evidence and qualification}\\
\midrule
\endfirsthead
\toprule
\textbf{Date} & \textbf{Building act or physical state} & \textbf{Evidence and qualification}\\
\midrule
\endhead
\bottomrule
\endfoot

1866 & Thirty families acquired and adapted a small brick church formerly used by Saint Stephen's Lutheran congregation at the northwest corner of Grove---now South Fifth---and Mineral Streets. & Gregory (1931), the National Register nomination, state inventory, and parish history agree on reuse. This was the first parish church, not the present Mitchell Street building.\\

Early 1870s & The parish acquired the Fifth-and-Mitchell site and financed a much larger church, remembered through a thirty-dollar family assessment and borrowing. & Parish and later press accounts supply the assessment. No subscription ledger, deed, loan instrument, or final project cost was examined.\\

July 1872 & Cornerstone laid for the present church, designed by Leonard A. Schmidtner. & HABS, Wisconsin AHI HI27246, Gregory, National Register nomination, and parish history converge. Later titles and biographical stories about Schmidtner are less secure than the building attribution.\\

1873 & A portion of the cream-brick, twin-towered church was dedicated and entered use. & Gregory and institutional histories use 1873 as dedication or completion; HABS shows that essential work continued later.\\

1870s & Early parish instruction acquired dedicated school space; Gregory dates a schoolhouse to 1878, while parish accounts give 1873 and other sources attach earlier dates to instruction. & The discrepancy likely separates teaching, organization, and building, but no checked record resolves the sequence.\\

1884--94 & Principal later church campaign: sanctuary built out, roof covered with copper, organ installed, and interior painted. & HABS WIS-159 is the controlling physical-history source. The range prevents the 1873 dedication from being mistaken for final completion.\\

1885 & Henry Messmer designed a substantial renovation valued in the state inventory at \$17,000. & Wisconsin AHI HI27246. The checked digital record does not itemize every component.\\

1894--95 & Interior brick walls were progressively faced with Wisconsin marble; four large bells were installed in the towers according to the detailed parish history. & Parish history gives 1894 and 1895 for marble stages and 1894 for the bells. A 2018 feature gives 1882 for the bells; no foundry record was checked.\\

1895--1913 & Gothic Revival entrance surround, revised tower-window archivolts, side buttresses with apostle statues, crossing flèche, marble altars, pulpit, and further interior elaboration. & HABS records the campaign but cannot divide all work by year. Later style labels describe layers, not one original scheme.\\

Late 19th--early 20th c. & Successive school, convent, rectory, and service buildings turned the site into a working campus. & Historic images, Gregory, HABS, parish history, and the National Register nomination establish a changing complex. A complete parcel-by-parcel chronology remains unavailable.\\

1926 & The principal four-story school front at 1669 South Fifth Street was built to plans by Herbst and Kuenzli, dated 16 February 1926; an older rear portion remained. & Wisconsin AHI HI118139 and the National Register nomination control the date and architect. The institutional phrase ``early 1930s'' is not used as the construction date.\\

Early 1930s & Church and parish buildings were redecorated or renovated during the Depression. & The parish history supplies the broad campaign. It does not identify every contract, room, or cost.\\

1932--47 & Saint Stanislaus High School opened in September 1932 with seventy-two freshmen; a 1946 report counted four hundred students from thirty-seven parishes; the state inventory records a 1947 change to the later Notre Dame High School name. & Contemporary \emph{Milwaukee Sentinel} report preserved by Marquette University and Wisconsin AHI HI118139. Enrollment figures are dated snapshots; educational expansion is not conflated with construction of the 1926 school front.\\

1960 & Twin towers tuckpointed; deteriorated dome structures removed and reconstructed in aluminum-covered steel, gilded with 23-carat gold leaf; crosses, clocks, and clock mechanisms renewed. & The ICKSP parish history gives the detailed 1960 sequence and September lifting of the towers. State and later records associate the golden domes more broadly with the early 1960s or 1966 centennial.\\

1962 & Fire damaged the left side altar and caused smoke damage throughout the church, delaying interior work. & Detailed institutional history. No fire-department report or insurance record was located.\\

1962/64 & Earlier stained-glass windows removed because of reported deterioration and security concerns; inch-thick dalle de verre glazing installed. & The City of Milwaukee survey says 1962; institutional history says 1964. No checked permit or contract resolves the date. Contractor and later press accounts document fragments from removed glass but not a complete survival.\\

1963--64 & Rectory comprehensively remade with modern interior arrangements and a concrete or limestone-facing treatment; ambulatory and chapel connected rectory and church. & Parish history, National Register nomination, and later reporting agree on a major project but differ over remodeling versus replacement. ``Remade'' preserves that uncertainty.\\

1965 & Roman travertine added to remaining interior walls; high altar dismantled or incorporated into a freestanding altar; pulpit moved; most communion rail removed; sanctuary extended into nave; baptismal font moved to new chapel. & Institutional history supplies the sequence. Later sources sometimes describe all changes generically as ``postconciliar,'' but the dated campaign had maintenance, safety, centennial, artistic, and liturgical dimensions.\\

Mid-1960s & Exterior mosaic shrine of Our Lady of Częstochowa installed as a parish-centennial and Polish-Christian-millennium memorial. & Institutional history. Its promotional size superlative was not independently measured.\\

2008 onward & Freestanding altar removed; high altar reconstructed; sanctuary location and floor recovered; marble communion rail and choir stalls introduced; historic furnishings and vestments returned to use. & ICKSP history and 2016 \emph{Catholic Herald} reporting. Components accrued in phases; the checked record does not support one single completion day or exact return to an untouched original state.\\

2011--12 & Organ pipes cleaned and regulated; action, bellows, and control system rebuilt; another console installed. & Detailed ICKSP parish history. No organ-builder contract or opus record was found in the checked set.\\

2013--14 & Priests' sacristy restored, including inlaid wood floor and cabinetry; 1960s chapel adapted as server sacristy, reserve-tabernacle location, and storage. & ICKSP parish history. The reported relic collection is a custodial use, not a building-ownership claim.\\

2014/15--21 & Historic fragments and other visual evidence informed a new stained-glass program; sanctuary, nave, rose, choir-loft, confessional, and baptistery windows installed in phases, replacing dalle de verre. & ICKSP history and restoration page, Conrad Schmitt portfolio, and 2016 reporting. Sources differ on discovery year; the institutional history dates completion to 2021.\\

2016 & Interior repainted from historical evidence; concealed marble cleaned and restored; rebuilt main altar consecrated during the parish's 150th-anniversary weekend. & ICKSP history and anniversary record; \emph{Catholic Herald} documents design research and phased fundraising. Consecration is a liturgical act distinct from the oratory's disputed erection date.\\

2017--18 & Choir loft redesigned; cabinetry, furniture, and new stairs added. & ICKSP parish history. This is distinct from the 2011--12 organ mechanism work.\\

2018--19 & Tower domes removed and rebuilt in copper; exterior brick cleaned and tuckpointed; roofs repaired; historically patterned wood doors installed; gardens, walks, and stairs renewed. & ICKSP history and Chrystine Elle Hanus's roofing-trade account, which dates the roofing project from December 2018 to December 2019. Copper replaced the failing 1960 aluminum-covered structures.\\

2019/20 & New energy-efficient chandeliers installed in designs intended to harmonize with the church; restored rectory rooms displayed as complete by 2020. & ICKSP prose dates chandelier installation to 2019 while a caption says 2020; the same history captions the restored rectory library 2020.\\

2023 & Scheduled maintenance found the striking ends of all four great bell clappers worn flat; replacement clappers restored tone and volume, and the originals were displayed in the Marian courtyard. & Saint Stanislaus parish bulletin, 29 July--6 August 2023. The replaceable striking parts were renewed; the nineteenth-century bells themselves remained in service.\\

March 2024 & Parish and oratory announced planned acquisition and renovation of two attached buildings across Fifth and Historic Mitchell for classrooms, meeting rooms, and parking, with an estimated twelve-to-eighteen-month construction period. & ICKSP announcement. It establishes intent, proposed scope, and an anonymous gift; later city records separately establish acquisition and approvals.\\

September 2024 snapshot & Archdiocesan Love One Another campaign page still identified grounds irrigation and exterior lighting as parish priorities. & Dated campaign record, not a total restoration budget. Its figures do not establish the aggregate cost of work since 2008.\\

17 October 2024 & St Stanislaus Properties LLC purchased the 501--505 West Historic Mitchell Street parcel for \$875,000. & City of Milwaukee assessor record, tax key 4620625000. The property record establishes the transaction and ownership name, not the internal relationship of the LLC to the parish or completion of renovations.\\

Jul.--Sept. 2025 & City bodies approved replacement brick facing within the historic district and a special use for a religious assembly hall. & Historic Preservation Commission record and Board of Zoning Appeals minutes. The special-use approval required permits and occupancy certificates within one year; approval is not proof that rehabilitation or occupancy was complete by the cutoff.\\
\end{longtable}
